\documentclass{article}

\usepackage{amsmath}
\usepackage{amssymb}
\usepackage{graphicx}
\usepackage{tikz}
\usepackage{pgfplots}
\pgfplotsset{compat=1.18}
\usepackage{booktabs}
\usepackage{xcolor}

\newenvironment{lesson-meta}{}{}
\newenvironment{item-analysis}{}{}
\newenvironment{model-discuss}{}{}
\newenvironment{concept-box}{}{}
\newenvironment{concept-summary}{}{}
\newenvironment{example}[1]{}{}
\newenvironment{tryit}[1]{}{}
\newenvironment{practice}[2]{}{}
\newenvironment{te-addendum}[2]{}{}
\newcommand{\answer}[1]{}
\newcommand{\placeholder}[2]{}
\newcommand{\uncertain}[2]{#1}
\newcommand{\te}[1]{}

\begin{document}

\begin{lesson-meta}
Lesson Code: 4-4

Title: Adding and Subtracting Rational Expressions

I CAN: I can find the sum or difference of rational expressions.

Objective: Students will be able to understand that rational expressions form a system analogous to the system of rational numbers and use that understanding to add and subtract rational expressions.

Essential Question: How do you rewrite rational expressions to find sums and differences?

Lesson Vocabulary: compound fraction.

Topic Vocabulary: algebraic fraction; numeric fraction.
\end{lesson-meta}

\begin{item-analysis}
\begin{tabular}{ccc}
\toprule
Example & Items & DOK \
\midrule
1 & 18--19 & 1 \
& 15 & 2 \
2 & 20--21 & 1 \
& 16 & 2 \
3 & 22--23 & 1 \
& 13, 17 & 2 \
4 & 24, 25, 35 & 1 \
& 12 & 2 \
5 & 26 & 2 \
& 31, 33 & 3 \
6 & 27--30 & 2 \
& 14, 32, 34 & 3 \
\bottomrule
\end{tabular}
\end{item-analysis}

\begin{model-discuss}
CRITIQUE & EXPLAIN

Teo and Evander find the following exercise in their homework:
[
\frac{1}{2}+\frac{1}{3}+\frac{1}{9}
]

A. Teo claims that a common denominator of the sum is (2+3+9=14). Evander claims that it is (2\cdot 3\cdot 9=54). Is either student correct? Explain why or why not.

B. Find the sum, explaining the method you use.

C. Construct Arguments Timothy states that the quickest way to find the sum of any two fractions with unlike denominators is to multiply their denominators to find a common denominator, and then rewrite each fraction with that denominator. Do you agree?

\placeholder{illustration}{Computer or tablet screen showing the expression (1/2+1/3+1/9).}

\answer{A. Only Evander is correct. A common denominator will have each other denominator as a factor. B. (\frac{17}{18}); converted each fraction to a fraction with denominator 18. C. No; It may be faster to use the least common denominator rather than multiplying the denominators together.}
\end{model-discuss}

\begin{te-addendum}{ModelDiscuss}{PurposefulQuestions}
Before: Implement Tasks that Promote Reasoning and Problem Solving.

Q: What do you notice about the expression in the homework exercise?
\answer{Three fractions are added. The numerator of each fraction is 1. Each fraction has a different denominator.}
\end{te-addendum}

\begin{te-addendum}{ModelDiscuss}{PurposefulQuestions}
During: Support Productive Struggle in Learning Mathematics.

Q: Compare common multiples and the least common multiple. How are they similar? How are they different?
\answer{The least common multiple is a common multiple. The common multiples of a group of numbers are multiples of each number in the group. A group of numbers has many common multiples, but only one least common multiple.}
\end{te-addendum}

\begin{te-addendum}{ModelDiscuss}{PurposefulQuestions}
For Early Finishers.

Q: Add (\frac{1}{8}) to the expression in the homework exercise. What are the least common denominator and the sum of the new expression?
\answer{72; (\frac{77}{72})}
\end{te-addendum}

\begin{te-addendum}{ModelDiscuss}{PurposefulQuestions}
After: Facilitate Meaningful Mathematical Discourse.

Q: Why would you want to use the least common multiple when writing fractions with a common denominator instead of any other common multiple?
\answer{The least common multiple gives you smaller numbers to compute.}

Q: How can there be more than one correct answer to this problem?
\answer{You can write the answer with different equivalent fractions. All of these answers to the homework exercise are equivalent: (\frac{17}{18}), (\frac{34}{36}), (\frac{51}{54}), and (\frac{68}{72}).}
\end{te-addendum}

\begin{te-addendum}{ModelDiscuss}{HabitsOfMind}
Look for Relationships For two fractions with denominators 10 and (x), when is (10x) the least common multiple? When is (10x) NOT the least common multiple?
\answer{When (x) and 10 have no common factors, (10x) is the least common multiple; when (x) and 10 have common factors, (10x) is not the least common multiple.}
\end{te-addendum}

\begin{concept-box}
Use Structure Compare addition of numerical and algebraic fractions:
[
\frac{1}{5}+\frac{3}{5}=\frac{1+3}{5}=\frac{4}{5}
]
In the same way,
[
\frac{x}{x+4}+\frac{5}{x+4}=\frac{x+5}{x+4}.
]
\end{concept-box}

\begin{example}{1}
Add Rational Expressions With Like Denominators

What is the sum?

A.
[
\frac{x}{x+4}+\frac{5}{x+4}
=\frac{x+5}{x+4}
]
The sum (\frac{x}{x+4}+\frac{5}{x+4}) is (\frac{x+5}{x+4}).

B.
[
\frac{2x+1}{x^2+3x}+\frac{3x-8}{x(x+3)}
=\frac{(2x+1)+(3x-8)}{x^2+3x}
]
[
=\frac{(2x+3x)+(1-8)}{x^2+3x}
=\frac{5x-7}{x^2+3x}
]
The sum (\frac{2x+1}{x^2+3x}+\frac{3x-8}{x(x+3)}) is (\frac{5x-7}{x^2+3x}).

\answer{A. (\frac{x+5}{x+4}). B. (\frac{5x-7}{x^2+3x}).}
\end{example}

\begin{tryit}{1}
Find the sum.

a. (\displaystyle \frac{10x-5}{2x+3}+\frac{8-4x}{2x+3})

b. (\displaystyle \frac{x-5}{x+5}+\frac{3x-21}{x+5})

\answer{a. (\frac{6x+3}{2x+3}). b. (\frac{4x-26}{x+5}) or (\frac{2(2x-13)}{x+5}).}
\end{tryit}

\begin{te-addendum}{1}{PurposefulQuestions}
Q: What do you notice about the denominators in Part A?
\answer{The denominators of the rational expressions are the same.}

Q: What do you notice about the denominators in Part B?
\answer{Even though the denominators are written differently, they are equivalent: (x^2+3x=x(x+3)).}
\end{te-addendum}

\begin{te-addendum}{1}{ElicitEvidence}
Q: Can you add the numerators of these rational expressions?
\answer{Yes; the denominators of the rational expressions are the same, so you can add the numerators.}
\end{te-addendum}

\begin{te-addendum}{1}{HabitsOfMind}
Make Sense and Persevere Explain why it does not make sense to add the denominators when adding rational numbers. Use numerical fractions to support your thinking.
\answer{When adding two fractions such as (\frac{1}{3}+\frac{2}{9}), add their numerators and keep the common denominator to get a sum of (\frac{5}{9}). You use the same process to add rational numbers.}
\end{te-addendum}

\begin{concept-box}
Use Structure The LCM of 6 and 15 is 30, not 90. The LCM does not contain the common factor 3 twice. Similarly, in Example 2B, the common factors ((x+3)) and ((x-5)) are not used twice.
\end{concept-box}

\begin{example}{2}
Identify the Least Common Multiple of Polynomials

How can you find the least common multiple (LCM) of polynomials?

A. ((x+2)^2,\ x^2+5x+6)

Factor each polynomial.
[
(x+2)^2=(x+2)(x+2)
]
[
x^2+5x+6=(x+2)(x+3)
]
The LCM is the product of the factors. Include the greatest number of any one factor in each polynomial.
[
\mathrm{LCM}:(x+2)(x+2)(x+3)\quad \text{or}\quad (x+2)^2(x+3)
]

B. (x^3-9x,\ x^2-2x-15,\ x^2-5x)

Factor each polynomial.
[
x^3-9x=x(x^2-9)=x(x+3)(x-3)
]
[
x^2-2x-15=(x+3)(x-5)
]
[
x^2-5x=x(x-5)
]
[
\mathrm{LCM}:x(x+3)(x-3)(x-5)
]

\answer{A. ((x+2)^2(x+3)). B. (x(x+3)(x-3)(x-5)).}
\end{example}

\begin{tryit}{2}
Find the LCM for each set of expressions.

a. (x^3+9x^2+27x+27,\ x^2-4x-21)

b. (10x^2-10y^2,\ 15x^2-30xy+15y^2,\ x^2+3xy+2y^2)

\answer{a. ((x+3)^3(x-7)). b. (30(x+y)(x+2y)(x-y)^2).}
\end{tryit}

\begin{te-addendum}{2}{PurposefulQuestions}
Q: How can you find the LCM of two rational expressions if the polynomials cannot be factored?
\answer{You multiply the two polynomials.}
\end{te-addendum}

\begin{te-addendum}{2}{ElicitEvidence}
Q: How could the LCM of two polynomials be one of the polynomials?
\answer{One of the polynomials is a factor of the other.}
\end{te-addendum}

\begin{te-addendum}{2}{RtISupport}
USE WITH EXAMPLE 2 Students may have difficulty finding the LCM of polynomials. Have students practice finding the LCM of integers.

1. Find the LCM of 24 and 42.
   \answer{168}

2. Find the LCM of 11 and 17.
   \answer{187}

3. Find the LCM of 15, 25, and 40.
   \answer{600}
   \end{te-addendum}

\begin{concept-box}
Make Sense and Persevere We must keep equivalent addends. So multiply each addend by a form of one. Choose the missing factors of the LCD.
\end{concept-box}

\begin{example}{3}
Add Rational Expressions With Unlike Denominators

What is the sum of (\frac{x+3}{x^2-1}) and (\frac{2}{x^2-3x+2})?

Follow a similar procedure to the one you use to add numerical fractions with unlike denominators.
[
\frac{x+3}{x^2-1}+\frac{2}{x^2-3x+2}
=\frac{x+3}{(x+1)(x-1)}+\frac{2}{(x-1)(x-2)}
]
[
=\frac{(x+3)(x-2)}{(x+1)(x-1)(x-2)}
+\frac{2(x+1)}{(x+1)(x-1)(x-2)}
]
[
=\frac{(x+3)(x-2)+2(x+1)}{(x+1)(x-1)(x-2)}
]
[
=\frac{(x^2+x-6)+(2x+2)}{(x+1)(x-1)(x-2)}
]
[
=\frac{x^2+3x-4}{(x+1)(x-1)(x-2)}
=\frac{(x+4)(x-1)}{(x+1)(x-1)(x-2)}
]
[
=\frac{x+4}{(x+1)(x-2)}
\quad \text{for } x\neq -1,\ 1,\ \text{and }2.
]

The sum of (\frac{x+3}{x^2-1}) and (\frac{2}{x^2-3x+2}) is (\frac{x+4}{(x+1)(x-2)}) for (x\neq -1,\ 1,) and (2).

\answer{(\frac{x+4}{(x+1)(x-2)}), for (x\neq -1,\ 1,) and (2).}
\end{example}

\begin{tryit}{3}
Find the sum.

a. (\displaystyle \frac{x+6}{x^2-4}+\frac{2}{x^2-5x+6})

b. (\displaystyle \frac{2x}{3x+4}+\frac{4x^2-11x-12}{6x^2+5x-4})

\answer{a. (\frac{x+7}{(x+2)(x-3)}), for (x\neq -2,\ 2,) and (3). b. (\frac{8x^2-13x-12}{(2x-1)(3x+4)}), for (x\neq -\frac{4}{3}) and (\frac{1}{2}).}
\end{tryit}

\begin{te-addendum}{3}{PurposefulQuestions}
Q: Why would you want to find the sum using LCM instead of any other common multiple?
\answer{Using the LCM will reduce the amount of simplifying you need to do after adding.}

Q: How is thinking about adding fractions with unlike denominators helpful when adding rational expressions with unlike denominators?
\answer{The process that you use to add fractions with unlike denominators is similar to the process for adding rational expressions with unlike denominators. Thinking about fractions may help simplify the process.}
\end{te-addendum}

\begin{te-addendum}{3}{ElicitEvidence}
Q: Why should you try to factor the expression after adding the numerators?
\answer{You factor the expression to see if there are any common factors you can divide out of the numerator and denominator.}
\end{te-addendum}

\begin{te-addendum}{3}{RtIExtend}
USE WITH EXAMPLE 3 Challenge students by having them add three rational expressions with unlike denominators.

Find the sum.

1. (\displaystyle \frac{3}{x^2-4x+3}+\frac{x+4}{x^2-1}+\frac{x}{x^2-2x-3})
   \answer{(\frac{2x^2+3x-9}{(x-1)(x-3)(x+1)})}

2. (\displaystyle \frac{x-5}{x^2+6x+8}+\frac{3-x}{x^2-3x-10}+\frac{4x+1}{x^2-x-20})
   \answer{(\frac{4x^2-2x+39}{(x+2)(x+4)(x-5)})}
   \end{te-addendum}

\begin{te-addendum}{4}{LearnTogether}
Additional Example 4 Help students understand adding and subtracting rational expressions with unlike denominators.

Q: What values are not in the domain when adding or subtracting rational expressions?
\answer{Any value that makes any denominator equal to 0 is not in the domain.}
\end{te-addendum}

\begin{concept-box}
Common Error When subtracting polynomials, remember to distribute (-1) when removing the parentheses.
\end{concept-box}

\begin{example}{4}
Subtract Rational Expressions

What is the difference between (\frac{x+1}{x^2-6x-16}) and (\frac{x+1}{x^2+6x+8})?
[
\frac{x+1}{x^2-6x-16}-\frac{x+1}{x^2+6x+8}
=\frac{x+1}{(x-8)(x+2)}-\frac{x+1}{(x+2)(x+4)}
]
The LCD is ((x-8)(x+2)(x+4)).
[
=\frac{(x+1)(x+4)}{(x-8)(x+2)(x+4)}
-\frac{(x-8)(x+1)}{(x-8)(x+2)(x+4)}
]
[
=\frac{(x^2+5x+4)-(x^2-7x-8)}{(x-8)(x+2)(x+4)}
]
[
=\frac{x^2+5x+4-x^2+7x+8}{(x-8)(x+2)(x+4)}
=\frac{12x+12}{(x-8)(x+2)(x+4)}
]
[
=\frac{12(x+1)}{(x-8)(x+2)(x+4)}
]
The difference between (\frac{x+1}{x^2-6x-16}) and (\frac{x+1}{x^2+6x+8}) is
[
\frac{12(x+1)}{(x-8)(x+2)(x+4)}
]
for (x\neq -4,\ -2,) and (8).

\answer{(\frac{12(x+1)}{(x-8)(x+2)(x+4)}), for (x\neq -4,\ -2,) and (8).}
\end{example}

\begin{tryit}{4}
Simplify.

a. (\displaystyle \frac{1}{3x}+\frac{1}{6x}-\frac{1}{x^2})

b. (\displaystyle \frac{3x-5}{x^2-25}-\frac{2}{x+5})

\answer{a. (\frac{x-2}{2x^2}), for (x\neq 0). b. (\frac{1}{x-5}), for (x\neq -5) and (5).}
\end{tryit}

\begin{te-addendum}{4}{PurposefulQuestions}
Q: How is subtracting rational expressions similar to adding rational expressions?
\answer{When adding or subtracting, you need to find the least common denominator and rewrite the rational expressions before performing the operation.}

Q: How is the Distributive Property used when subtracting rational expressions?
\answer{The Distributive Property is used when simplifying the numerator by subtracting polynomials. You multiply each term of the polynomial by (-1).}
\end{te-addendum}

\begin{te-addendum}{4}{CommonError}
Try It 4b When subtracting rational expressions, students may forget to distribute the negative sign. Have students relate subtracting rational expressions to subtracting polynomials, reminding them that they need to multiply each term by (-1) before removing the parentheses and combining like terms in the numerator.
\end{te-addendum}

\begin{te-addendum}{4}{HabitsOfMind}
Communicate Precisely How does finding the LCM of two or more polynomials help you to add and subtract rational expressions?
\answer{Once you know the LCM of polynomials in the denominators, you can multiply and divide each rational expression by the expression required to make each denominator the LCM. Then you can add or subtract the rational expressions.}
\end{te-addendum}

\begin{concept-box}
Use Appropriate Tools Use a table to organize information and help create an accurate model of the situation.
\end{concept-box}

\begin{example}{5}
Use Rational Expressions to Solve Problems

Jorge drives his car to the mechanic, then he takes the commuter rail train back to his neighborhood. The average speed for the 10-mile trip is 15 miles per hour faster on the train. Find an expression for Jorge's total travel time. If he drove 30 mph, how long did this take?

\placeholder{illustration}{Commuter rail train and car; train label "Speed = r + 15" and car label "Speed = r".}

[
\begin{array}{lccc}
\toprule
& \text{Distance} & \text{Rate} & \text{Time} \
\midrule
\text{Car} & 10 & r & \frac{10}{r} \
\text{Commuter Rail} & 10 & r+15 & \frac{10}{r+15} \
\bottomrule
\end{array}
]

Remember: distance (=) rate (\cdot) time, so time (=\frac{\text{distance}}{\text{rate}}).

Total time for the trip:
[
\frac{10}{r}+\frac{10}{r+15}
=\frac{10(r+15)}{r(r+15)}+\frac{10r}{r(r+15)}
]
[
=\frac{10r+150+10r}{r(r+15)}
=\frac{20r+150}{r(r+15)}
]

At a driving rate of 30 mph, you can find the total time.
[
\frac{20r+150}{r(r+15)}
=\frac{20(30)+150}{30(30+15)}
=\frac{750}{1350}
=\frac{5}{9}
]
The expression for Jorge's total travel time is (\frac{20r+150}{r(r+15)}). The total time is (\frac{5}{9}) h, or about 33 min.

\answer{(\frac{20r+150}{r(r+15)}); (\frac{5}{9}) h, or about 33 min.}
\end{example}

\begin{tryit}{5}
Jorge drives to visit his friends. His 20 mile trip to their house is congested with traffic, but it has cleared by the time he travels home. He can travel 5 miles per hour faster. Write an expression that represents Jorge's total travel time.

\answer{(\frac{40r+100}{r(r+5)}) or (\frac{20(2r+5)}{r(r+5)}).}
\end{tryit}

\begin{te-addendum}{5}{PurposefulQuestions}
Q: How do you know that you are supposed to add the two rational expressions?
\answer{The problem asks you to find an expression for the total time. Each expression represents the time for part of the trip. You need to add the times to find the total time.}

Q: How is the table a useful tool when solving this problem?
\answer{The table organizes all of the information in the problem statement and the picture in one place. From the table, it is easy to see the expressions needed to calculate the total time.}
\end{te-addendum}

\begin{te-addendum}{5}{HabitsOfMind}
Construct Arguments Does Jorge spend more time traveling to his friend's house, or back home? Support your answer algebraically.
\answer{Traveling to his friend's house. You can express the time traveled to Jorge's friend's house by (\frac{20}{r}) and the time home by (\frac{20}{r+5}). Since (\frac{20}{r+5}<\frac{20}{r}), Jorge spends less time traveling home.}
\end{te-addendum}

\begin{te-addendum}{5}{LearnTogether}
Additional Example 5 Help students transition to solving a rate problem for a three-part trip.

Q: Can you determine the time of the trip with the information given? Explain.
\answer{No; you need to know the speed that Franklin walked.}
\end{te-addendum}

\begin{concept-box}
Vocabulary A compound fraction contains a rational expression in its numerator, denominator, or both.
\end{concept-box}

\begin{example}{6}
Simplify a Compound Fraction

Write a simpler form of the compound fraction
[
\frac{\frac{1}{x}+\frac{2}{x+1}}{\frac{1}{y}}.
]

Find the Least Common Multiple (LCM) of all the denominators.
[
\frac{\frac{1}{x}+\frac{2}{x+1}}{\frac{1}{y}}
=============================================

\frac{\left(\frac{1}{x}+\frac{2}{x+1}\right)\cdot xy(x+1)}
{\left(\frac{1}{y}\right)\cdot xy(x+1)}
]
[
=

\frac{\frac{1}{x}\cdot xy(x+1)+\frac{2}{x+1}\cdot xy(x+1)}
{\frac{1}{y}\cdot xy(x+1)}
]
[
=

# \frac{(x+1)y+2xy}{x(x+1)}

# \frac{xy+y+2xy}{x(x+1)}

\frac{(3x+1)y}{x(x+1)}
]
So,
[
\frac{\frac{1}{x}+\frac{2}{x+1}}{\frac{1}{y}}
]
is equivalent to
[
\frac{(3x+1)y}{x(x+1)}
]
when (x\neq -1,\ 0) and (y\neq 0).

\answer{(\frac{(3x+1)y}{x(x+1)}), when (x\neq -1,\ 0) and (y\neq 0).}
\end{example}

\begin{tryit}{6}
Simplify each compound fraction.

a. (\displaystyle \frac{\frac{1}{x-1}}{\frac{x+1}{3}+\frac{4}{x-1}})

b. (\displaystyle \frac{2-\frac{1}{x}}{x+\frac{2}{x}})

\answer{a. (\frac{3}{x^2+11}), for (x\neq 1). b. (\frac{2x-1}{x^2+2}), for (x\neq 0).}
\end{tryit}

\begin{te-addendum}{6}{PurposefulQuestions}
Q: What do you notice about a compound fraction that is different from other fractions?
\answer{The numerator or denominator of a compound fraction contains a fraction.}

Q: Is there another method you could use to write a simpler form of the compound fraction?
\answer{You can multiply the numerator by the reciprocal of the denominator, then write the rational expression as a single fraction.}
\end{te-addendum}

\begin{te-addendum}{6}{ElicitEvidence}
Q: What first step would you take to rewrite the compound fraction?
\answer{The first step is to multiply the numerator and denominator by the LCD of all rational expressions in the numerator and denominator. This eliminates the fractions from the numerator and denominator and makes the compound fraction easier to simplify.}
\end{te-addendum}

\begin{te-addendum}{6}{HabitsOfMind}
Reason Edwin multiplied the top and bottom of the fraction in Try It 6a by (\frac{3}{x+1}+\frac{x-1}{4}). Will this technique work to simplify the compound fraction? Explain.
\answer{No; this will not simplify the denominator because the product of (\frac{x+1}{3}+\frac{4}{x-1}) and (\frac{3}{x+1}+\frac{x-1}{4}) gives 2 plus additional terms.}
\end{te-addendum}

\begin{te-addendum}{6}{ELLAddendum}
READING Beginning Help students understand the steps in the method by previewing the math vocabulary before reading the example. Show the students the following terms and their definitions: LCD, numerator, denominator, and reciprocal. Give students a copy of the math steps in the example. Have them read the example and label the terms as they are reading.

Q: How is the LCD different from the denominator in the compound fraction?
\answer{The LCD is the product of all of the denominators that appear in the compound fraction (x), (y), and (x+1). The denominator of the compound fraction is (\frac{1}{y}).}

SPEAKING Developing As you review the steps in simplifying a compound fraction, discuss the terms eliminate and reciprocal. Have students take turns explaining the process in each method.

Q: When multiplying the numerator and denominator of the compound fraction by the LCD, what is being eliminated?
\answer{Multiplying the numerator and denominator by the LCD ``eliminates'' fractions. The compound fraction is rewritten with a numerator and denominator with no fractions.}

Q: What property is helpful in eliminating the fractions?
\answer{The Distributive Property is helpful for multiplying each addend in brackets by the LCD.}

LISTENING Expanding Have students work in groups. Give each group cards labeled with numbers, operations, and fractions. Read aloud the problem statement from the example about a compound fraction. As they listen to the description, have students arrange the cards to create a compound fraction.

Q: Does a compound fraction have to include more than one fraction?
\answer{No; the description says one or more fractions.}

Q: Does a compound fraction have to have fractions in both the numerator and denominator?
\answer{No; the definition says and/or, which means a compound fraction may have a fraction in either the numerator or denominator, or it may have fractions in both.}
\end{te-addendum}

\begin{concept-summary}
CONCEPT SUMMARY Find Sums and Differences of Rational Expressions

Words: To add or subtract rational expressions with common denominators, add the numerators and keep the denominator the same.
[
\frac{1}{5}+\frac{3}{5}=\frac{1+3}{5}=\frac{4}{5}
]
[
\frac{x}{x+4}+\frac{5}{x+4}=\frac{x+5}{x+4}
]

Words: To add or subtract rational expressions with different denominators, rewrite each expression so that its denominator is the LCD, then add or subtract the numerators.
[
\frac{1}{6}+\frac{1}{15}
=\frac{1}{2\cdot 3}+\frac{1}{3\cdot 5}
=\frac{1\cdot 5+1\cdot 2}{2\cdot 3\cdot 5}
=\frac{7}{30}
]
[
\frac{x+3}{x^2-1}+\frac{2}{x^2-3x+2}
====================================

\frac{x+3}{(x+1)(x-1)}+\frac{2}{(x-1)(x-2)}
]
Rewrite the rational expressions using the LCD:
[
\frac{(x+3)(x-2)}{(x+1)(x-1)(x-2)}
+
\frac{2(x+1)}{(x+1)(x-1)(x-2)}.
]

Do You UNDERSTAND?

1. ESSENTIAL QUESTION How do you rewrite rational expressions to find sums and differences?
   \answer{Find a common denominator, add or subtract the numerators, and simplify, eliminating common factor(s) shared by both the numerator and denominator of the sum or difference.}

2. Vocabulary In your own words, define compound fraction and provide an example of one.
   \answer{A compound fraction is a fraction that has one or more fractions in the numerator and/or denominator. Example: (\frac{\frac{1}{2x}}{x+y}).}

3. Error Analysis A student added the rational expressions as follows:
   [
   \frac{5x}{x+7}+\frac{7}{x}
   =
   \frac{5x}{x+7}+\frac{7(7)}{x+7}
   =
   \frac{5x+49}{x+7}.
   ]
   Describe and correct the error the student made.
   \answer{The student incorrectly used the polynomial expression (x+7) as the common denominator, when (x(x+7)) should have been used. The correct method is (\frac{5x}{x+7}+\frac{7}{x}=\frac{5x(x)}{x(x+7)}+\frac{7(x+7)}{x(x+7)}=\frac{5x^2+7x+49}{x(x+7)}).}

4. Construct Arguments Explain why, when stating the domain of a sum or difference of rational expressions, not only should the simplified sum or difference be considered but the original expression should also be considered.
   \answer{The simplified sum or difference may not include some values of the domain that were eliminated during the simplification process.}

5. Make Sense and Persevere In adding or subtracting rational expressions, why is the L in LCD significant?
   \answer{LCD stands for least common denominator. The sum or difference will not be in simplest form if you multiply the numerator and denominator of each expression by more factors than necessary.}

Do You KNOW HOW?

6. Find the sum (\displaystyle \frac{3}{x+1}+\frac{11}{x+1}).
   \answer{(\frac{14}{x+1})}

Find the LCM of the polynomials.

7. (x^2-y^2) and (x^2-2xy+y^2)
   \answer{((x+y)(x-y)^2)}

8. (5x^3y) and (15x^2y^2)
   \answer{(15x^3y^2)}

Find the sum or difference.

9. (\displaystyle \frac{3x}{4y^2}-\frac{y}{10x})
   \answer{(\frac{15x^2-2y^3}{20xy^2})}

10. (\displaystyle \frac{9y+2}{3y^2-2y-8}+\frac{7}{3y^2+y-4})
    \answer{(\frac{3y-4}{(y-2)(y-1)})}

11. Find the perimeter of the quadrilateral in simplest form.
    [
    \begin{tikzpicture}[scale=0.75]
    \draw (0,0) -- (1.1,1.8) -- (3.6,2.2) -- (5.2,0.8) -- cycle;
    \node at (1.8,2.35) {(x)};
    \node at (4.75,1.65) {(\frac{2}{3x})};
    \node at (2.5,-0.25) {(2x)};
    \node at (0.15,1.0) {(\frac{1}{x})};
    \end{tikzpicture}
    ]
    \answer{(\frac{9x^2+5}{3x}) units}
    \end{concept-summary}

\begin{te-addendum}{ConceptSummary}{PurposefulQuestions}
Q: Are the steps for subtracting two rational expressions the same as the steps for adding two rational expressions?
\answer{When subtracting two rational expressions, you must remember to distribute the negative sign to each term in the numerator of the expression being subtracted.}
\end{te-addendum}

\begin{te-addendum}{ConceptSummary}{CommonError}
Exercise 7 Students may incorrectly find the LCM by multiplying the two polynomials. Have students factor each polynomial and identify any repeated factors before they multiply to find the LCM.
\end{te-addendum}

\begin{practice}{12}{2}
Generalize Explain how addition and subtraction of rational expressions is similar to and different from addition and subtraction of rational numbers.
\answer{The same basic procedures are followed for both. First, find a common denominator, then rewrite the expressions using the common denominator, perform the addition and/or subtraction, and then simplify.}
\end{practice}

\begin{practice}{13}{2}
Error Analysis Describe and correct the error a student made in adding the rational expressions.
[
\frac{1}{x^2+3x+2}+\frac{x^2+4x}{4x+8}
======================================

\frac{1}{(x+1)(x+2)}+\frac{x(x+4)}{4(x+2)}
]
[
=

\frac{4}{4(x+1)(x+2)}+\frac{x(x+4)}{4(x+1)(x+2)}
]
[
=

# \frac{4+x^2+4x}{4(x+1)(x+2)}

\frac{x^2+4x+4}{4(x+1)(x+2)}
]
[
=

# \frac{(x+2)(x+2)}{4(x+1)(x+2)}

\frac{x+2}{4(x+1)}
]
\answer{When rewriting the second expression with a common denominator, the student should have had (x(x+1)(x+4)) in the numerator rather than (x(x+4)).}
\end{practice}

\begin{practice}{14}{3}
Use Structure Find the slope of the line that passes through the points shown. Express in simplest form.
[
\begin{tikzpicture}
\begin{axis}[
axis lines=middle,
xmin=-4.5,xmax=4.5,
ymin=-2.5,ymax=4.5,
xlabel={(x)},
ylabel={(y)},
xtick={-4,-2,2,4},
ytick={-2,2,4},
width=7cm,
height=5.5cm,
grid=none,
clip=false
]
\addplot[domain=-4:4] {x};
\addplot[only marks, mark=*] coordinates {(1,1) (-1,-1)};
\node[anchor=west] at (axis cs:1,1) {(\left(\frac{1}{a},\frac{1}{b}\right))};
\node[anchor=east] at (axis cs:-1,-1) {(\left(-\frac{1}{b},-\frac{1}{a}\right))};
\end{axis}
\end{tikzpicture}
]
\answer{1}
\end{practice}

\begin{practice}{15}{2}
Reason For what values of (x) is the sum of (\frac{x-5y}{x+y}) and (\frac{x+y}{x+y}) undefined? Explain.
\answer{When (x=-y), the denominator is zero, which is undefined.}
\end{practice}

\begin{practice}{16}{2}
Look for Relationships Find the LCM of (3x^2+7x+2) and (9x+3).
\answer{(3(3x+1)(x+2))}
\end{practice}

\begin{practice}{17}{2}
Higher Order Thinking Why are rational expressions closed under addition?
\answer{You can find the sum of two rational expressions (\frac{p(x)}{q(x)}+\frac{r(x)}{t(x)}) by finding the sum of equivalent expressions, (\frac{p(x)\cdot t(x)}{q(x)\cdot t(x)}+\frac{r(x)\cdot q(x)}{t(x)\cdot q(x)}=\frac{p(x)\cdot t(x)+r(x)\cdot q(x)}{q(x)\cdot t(x)}). Since the numerator and denominator are still polynomials, the sum is a rational expression. So rational expressions are closed under addition.}
\end{practice}

\begin{practice}{18}{1}
Find the sum.
[
\frac{4x}{x+7}+\frac{9}{x+7}
]
\answer{(\frac{4x+9}{x+7})}
\end{practice}

\begin{practice}{19}{1}
Find the sum.
[
\frac{3y-1}{y^2+4y}+\frac{9y+6}{y(y+4)}
]
\answer{(\frac{12y+5}{y^2+4})}
\end{practice}

\begin{practice}{20}{1}
Find the LCM for each group of expressions.
[
x^2-7x+6,\quad x^2-5x-6
]
\answer{((x-6)(x+1)(x-1))}
\end{practice}

\begin{practice}{21}{1}
Find the LCM for each group of expressions.
[
y^2+2y-24,\quad y^2-16,\quad 2y
]
\answer{(2y(y+6)(y+4)(y-4))}
\end{practice}

\begin{practice}{22}{1}
Find the sum.
[
\frac{6x}{x^2-8x}+\frac{4}{2x-16}
]
\answer{(\frac{8}{x-8})}
\end{practice}

\begin{practice}{23}{1}
Find the sum.
[
\frac{3y}{2y^2-y}+\frac{2}{2y}
]
\answer{(\frac{5y-1}{y(2y-1)})}
\end{practice}

\begin{practice}{24}{1}
Find the difference.
[
\frac{4x}{x^2-1}-\frac{4}{x-1}
]
\answer{(-\frac{4}{(x+1)(x-1)})}
\end{practice}

\begin{practice}{25}{1}
Find the difference.
[
\frac{y-1}{3y+15}-\frac{y+3}{5y+25}
]
\answer{(\frac{2y-14}{15(y+5)}) or (\frac{2(y-7)}{15(y+5)})}
\end{practice}

\begin{practice}{26}{2}
On Saturday morning, Ahmed decided to take a bike ride from one end of the 15-mile bike trail to the other end of the bike trail and back. His average speed the first half of the ride was 2 mph faster than his speed the second half. Find an expression for Ahmed's total travel time. If his average speed for the first half of the ride was 12 mph, how long was Ahmed's bike ride?

\placeholder{illustration}{Bike trail with cyclist, sign "Trail length 15 miles", one label "Speed x+2" and another label "Speed x".}

\answer{(\frac{30(x+1)}{x(x+2)}); 2.75 hours}
\end{practice}

\begin{practice}{27}{2}
Write each expression in its simplest form.
[
\frac{1+\frac{1}{x}}{x-\frac{1}{x}}
]
\answer{(\frac{1}{x-1})}
\end{practice}

\begin{practice}{28}{2}
Write each expression in its simplest form.
[
\frac{\frac{3}{y}+\frac{7}{x}}{\frac{1}{y}-\frac{2}{x}}
]
\answer{(\frac{3x+7y}{x-2y})}
\end{practice}

\begin{practice}{29}{2}
Write each expression in its simplest form.
[
\frac{\frac{1}{a}+\frac{1}{b}}{\frac{a^2-b^2}{ab}}
]
\answer{(\frac{1}{a-b})}
\end{practice}

\begin{practice}{30}{2}
Write each expression in its simplest form.
[
\frac{\frac{z^2-z-12}{z^2-2z-15}}{\frac{z^2+8z+12}{z^2-5z-14}}
]
\answer{(\frac{(z-4)(z-7)}{(z-5)(z+6)})}
\end{practice}

\begin{practice}{31}{3}
Use Structure Aisha paddles a kayak 5 miles downstream at a rate 3 mph faster than the river's current. She then travels 4 miles back upstream at a rate 1 mph slower than the river's current. Hint: Let (x) represent the rate of the river current.

\placeholder{illustration}{Kayak on river with arrows and labels "Kayak's downstream speed is x+3 mph" and "Water current's speed is x mph".}

a. Write and simplify an expression to represent the total time it takes Aisha to paddle the kayak 5 miles downstream and 4 miles upstream.

b. If the rate of the river current, (x), is 2 mph, how long was Aisha's entire kayak trip?

\answer{a. (\frac{9x+7}{(x+3)(x-1)}). b. 5 h.}
\end{practice}

\begin{practice}{32}{3}
Model With Mathematics The resistance of electrical circuits A and B are shown. Find the ratio of the resistance of Circuit A to Circuit B.

[
\begin{tikzpicture}[scale=0.85]
\draw (0,0) rectangle (3,1.6);
\node at (1.5,0.8) {A};
\node[left] at (0,0.8) {(\frac{x^2-25}{x-4}\ \Omega)};
\draw (3,0.55) -- (3.3,0.55);
\draw (3,1.05) -- (3.3,1.05);
\draw (4.2,0.2) rectangle (6.2,1.4);
\node at (5.2,0.8) {B};
\node[right] at (6.2,0.8) {(\frac{x+5}{x^2-16}\ \Omega)};
\draw (4.2,0.55) -- (3.9,0.55);
\draw (4.2,1.05) -- (3.9,1.05);
\end{tikzpicture}
]
\answer{((x-5)(x+4):1) or (x^2-x-20:1)}
\end{practice}

\begin{practice}{33}{3}
Reason The Bello family drives 180 miles (round trip) to a professional basketball game. On the way to the game, their average speed is approximately 8 mph faster than their speed on the return trip home.

a. Let (x) represent their average speed on the way home. Write and simplify an expression to represent the total time it took them to drive to and from the game.

b. If their average speed going to the game was 62 mph, how long did it take them to drive to the game and back?

\answer{a. (\frac{180x+720}{x(x+8)}) or (\frac{180(x+4)}{x(x+8)}). b. (\approx 3.1) h.}
\end{practice}

\begin{practice}{34}{3}
Which of the following compound fractions simplifies to (\frac{x+1}{x-3})? Select all that apply.

[
\begin{array}{ll}
\text{A. } \displaystyle
\frac{\frac{x^2+5x+4}{x^2+2x-8}}
{\frac{x^2-4x+3}{x^2-3x+2}}
&
\text{B. } \displaystyle
\frac{\frac{x^2-1}{x^2-4}}
{\frac{x^2+x-7}{x^2+5x+6}}
[3em]
\text{C. } \displaystyle
\frac{\frac{x^2+3x-10}{x^2-16}}
{\frac{x^2-4x-5}{x^2-1}}
&
\text{D. } \displaystyle
\frac{\frac{x^2+3x-10}{x^2-5x+6}}
{\frac{x^2-25}{x^2-4x-5}}
\end{array}
]
\answer{A and D}
\end{practice}

\begin{practice}{35}{1}
SAT/ACT What is the difference between (\frac{x}{9}) and (\frac{x-y}{6})?

A. (\displaystyle \frac{5x-y}{18})

B. (\displaystyle \frac{5x+y}{18})

C. (\displaystyle \frac{-x+3y}{18})

D. (\displaystyle \frac{-x-3y}{18})

\answer{C. (\frac{-x+3y}{18})}
\end{practice}

\begin{practice}{36}{\uncertain{3}{not listed in Item Analysis table}}
Performance Task The lens equation
[
\frac{1}{f}=\frac{1}{d_i}+\frac{1}{d_o}
]
represents the relationship between (f), the focal length of a camera lens, (d_i), the distance from the lens to the film, and (d_o), the distance from the lens to the object.

[
\begin{tikzpicture}[scale=0.75]
\draw[->] (0,0) -- (0,1.4);
\node[above] at (0,1.4) {Object};
\draw (4,0) ellipse (0.25 and 1.4);
\node[above] at (4,1.6) {Lens};
\draw (7,0) -- (7,2.1);
\node[above] at (7,2.1) {Image screen};
\draw[rounded corners] (6.4,-0.4) rectangle (8.1,2.4);
\draw (8.1,0.2) circle (0.35);
\draw[dashed] (0,1.2) -- (4,0.5) -- (7,1.2);
\draw[dashed] (0,1.2) -- (4,-0.1) -- (7,-0.8);
\draw[<->] (0,-0.8) -- (4,-0.8);
\node[below] at (2,-0.8) {(d_o)};
\draw[<->] (4,-0.8) -- (7,-0.8);
\node[below] at (5.5,-0.8) {(d_i)};
\draw[<->] (2.7,-0.25) -- (4,-0.25);
\node[below] at (3.35,-0.25) {(f)};
\draw[<->] (4,-0.25) -- (5.3,-0.25);
\node[below] at (4.65,-0.25) {(f)};
\end{tikzpicture}
]

Part A Find the focal length of a camera lens if an object that is 12 cm from a camera lens is in focus on the film when the lens is 6 cm from the film.

Part B Suppose the focal length of another camera lens is 3 inches, and the object to be photographed is 5 feet away. What distance (to the nearest tenth inch) should the lens be from the image screen?

\answer{Part A (f=4) cm. Part B (x\approx 3.2) inches.}
\end{practice}

\end{document}
